% Supplement accompanying the Wireless Networks formatted manuscript.
\documentclass[10pt,a4paper]{article}
\usepackage[top=25.4mm,bottom=25.4mm,left=31.75mm,right=31.75mm]{geometry}
\usepackage{iftex}
\ifPDFTeX
  \usepackage[T1]{fontenc}
  \usepackage{newtxtext}
\else
  \usepackage{fontspec}
  \IfFontExistsTF{Times New Roman}{\setmainfont{Times New Roman}}{\setmainfont{Tinos}}
\fi
\usepackage{booktabs}
\usepackage{adjustbox}
\usepackage[hyphens]{url}
\usepackage{amsmath}
\usepackage{setspace}
\setstretch{1.15}
\setlength{\emergencystretch}{2em}
\begin{document}
\begin{center}
{\large\bfseries Online Resource 1}\\[0.65em]
{\large False-alarm budget audit and duration sensitivity}\\[0.5em]
Evaluating False-Alarm Budgets for Fragmentation Warnings in Simulated Flying Ad Hoc Networks\\[0.5em]
Supporting material --- reproducible research draft
\end{center}

\subsection*{A. Development block by mobility}
\noindent Table A uses the fixed threshold of 0.15 and full simulated exposure. Each cell contains three trajectory groups and 12 runs. Four radio/graph conditions share each underlying trajectory. The denominator is full simulated time, not alert-eligible time. Meets budget denotes at most two false alarms per simulated minute.

\begin{center}
\small
\setlength{\tabcolsep}{4pt}
\begin{adjustbox}{max width=\linewidth}\begin{tabular}{llrrrrrrrr}
\toprule
Block & Mobility & Groups & Runs & Events & Matched & False alarms & Min & False/min & Meets \\
\midrule
% BEGIN VERIFIED R9.10 TABLE
2 & Random waypoint (RWP) & 3 & 12 & 79 & 2 & 2 & 4 & 0.50 & Yes \\
2 & Gauss--Markov (GM) & 3 & 12 & 56 & 25 & 10 & 4 & 2.50 & No \\
2 & Mission-based & 3 & 12 & 36 & 10 & 7 & 4 & 1.75 & Yes \\
3 & Random waypoint (RWP) & 3 & 12 & 27 & 0 & 0 & 4 & 0.00 & Yes \\
3 & Gauss--Markov (GM) & 3 & 12 & 68 & 15 & 6 & 4 & 1.50 & Yes \\
3 & Mission-based & 3 & 12 & 76 & 3 & 2 & 4 & 0.50 & Yes \\
4 & Random waypoint (RWP) & 3 & 12 & 100 & 4 & 1 & 4 & 0.25 & Yes \\
4 & Gauss--Markov (GM) & 3 & 12 & 100 & 13 & 9 & 4 & 2.25 & No \\
4 & Mission-based & 3 & 12 & 74 & 0 & 0 & 4 & 0.00 & Yes \\
5 & Random waypoint (RWP) & 3 & 12 & 48 & 13 & 9 & 4 & 2.25 & No \\
5 & Gauss--Markov (GM) & 3 & 12 & 123 & 34 & 10 & 4 & 2.50 & No \\
5 & Mission-based & 3 & 12 & 98 & 0 & 0 & 4 & 0.00 & Yes \\
6 & Random waypoint (RWP) & 3 & 12 & 51 & 5 & 6 & 4 & 1.50 & Yes \\
6 & Gauss--Markov (GM) & 3 & 12 & 118 & 30 & 14 & 4 & 3.50 & No \\
6 & Mission-based & 3 & 12 & 60 & 29 & 12 & 4 & 3.00 & No \\
% END VERIFIED R9.10 TABLE
\bottomrule
\end{tabular}\end{adjustbox}
\end{center}

\noindent Six cells exceed the budget. The three mobility regimes are repeated within each block. All five blocks are development data and were evaluated without fitting a new calibrator or tuning the warning threshold. The source data and independent verifier accompany the manuscript in the reproducibility package.

\subsection*{B. Mobility by node count}
\noindent Table B uses the same fixed threshold of 0.15. Each row covers five trajectory groups, one group in each development block, with four paired radio/graph conditions per group. Thus, each row represents 20 runs and \(20/3\) minutes of full simulated time. A dash denotes undefined precision when no alarms were emitted. Bold indicates more than two false alarms per simulated minute.

\begin{center}
\small
\setlength{\tabcolsep}{3.2pt}
\begin{adjustbox}{max width=\linewidth}\begin{tabular}{lrrrrrrrr}
\toprule
Mobility & Nodes & Runs & Events & Matched & False alarms & Recall (\%) & Precision (\%) & False/min \\
\midrule
% BEGIN VERIFIED R9.10 BY MOBILITY AND NODE COUNT
Random waypoint (RWP) & 10 & 20 & 181 & 24 & 17 & 13.3 & 58.5 & \textbf{2.55} \\
Random waypoint (RWP) & 20 & 20 & 113 & 0 & 1 & 0.0 & 0.0 & 0.15 \\
Random waypoint (RWP) & 30 & 20 & 11 & 0 & 0 & 0.0 & -- & 0.00 \\
Gauss--Markov (GM) & 10 & 20 & 205 & 92 & 39 & 44.9 & 70.2 & \textbf{5.85} \\
Gauss--Markov (GM) & 20 & 20 & 205 & 25 & 10 & 12.2 & 71.4 & 1.50 \\
Gauss--Markov (GM) & 30 & 20 & 55 & 0 & 0 & 0.0 & -- & 0.00 \\
Mission-based & 10 & 20 & 213 & 42 & 21 & 19.7 & 66.7 & \textbf{3.15} \\
Mission-based & 20 & 20 & 115 & 0 & 0 & 0.0 & -- & 0.00 \\
Mission-based & 30 & 20 & 16 & 0 & 0 & 0.0 & -- & 0.00 \\
% END VERIFIED R9.10 BY MOBILITY AND NODE COUNT
\bottomrule
\end{tabular}\end{adjustbox}
\end{center}

\noindent In the 30-node runs, all 11,285 connected decision scores were below the frozen cutoff of 0.15 (maximum 0.1160); no alarms were emitted, although 82 fragmentation events were observed. The mobility and node-count summaries describe the same 180 development runs. They are not new experiments or estimates of a causal mobility or network-size effect.

\medskip
\noindent Author and correspondence metadata are withheld in this review draft. Complete the corresponding fields for the submission version as instructed by the journal.
\newpage
\begin{center}
{\large\bfseries Online Resource 1 (continued)}
\end{center}

\subsection*{C. Fragmentation episode lengths}
\noindent Episode length is the number of consecutive sampled frames with a disconnected graph from an observed onset. Frames are sampled at 0.1-second intervals. If the run ends before recovery is observed, the length is a lower bound. An episode already in progress at the first frame does not have an observed onset.

\begin{center}
\small
\begin{adjustbox}{max width=\linewidth}\begin{tabular}{lrrr}
\toprule
Observed frames & Recovery observed & Right-censored & Total \\
\midrule
% BEGIN VERIFIED R9.10 EPISODE LENGTH DISTRIBUTION
1 & 592 & 7 & 599 \\
2 & 197 & 1 & 198 \\
3--5 & 183 & 1 & 184 \\
6--10 & 63 & 3 & 66 \\
11 or more & 61 & 6 & 67 \\
% END VERIFIED R9.10 EPISODE LENGTH DISTRIBUTION
\bottomrule
\end{tabular}\end{adjustbox}
\end{center}

\subsection*{D. Minimum-duration sensitivity}
\noindent The original one-frame definition is the reference. For each alternative, the same 271 alarms are rematched to confirmed onsets within the original six-frame window over 60 simulated minutes. An unconfirmed onset lacks enough observed frames to classify its duration. An unresolved alarm has no confirmed match but an unconfirmed onset in its window. Other unmatched alarms are false for the duration-specific target. Recall is matched divided by confirmed qualifying onsets. Lead is the median time from a matched alarm to onset, not to duration confirmation.

\begin{center}
\small
\setlength{\tabcolsep}{3.8pt}
\begin{adjustbox}{max width=\linewidth}\begin{tabular}{rrrrrrrrr}
\toprule
Min. frames & Confirmed & Unconfirmed & Matched & Recall (\%) & False & Unresolved & False/min & Lead (s) \\
\midrule
% BEGIN VERIFIED R9.10 BY MINIMUM EPISODE LENGTH
1 & 1114 & 0 & 183 & 16.4 & 88 & 0 & 1.47 & 0.20 \\
2 & 515 & 7 & 136 & 26.4 & 135 & 0 & 2.25 & 0.20 \\
3 & 317 & 8 & 105 & 33.1 & 166 & 0 & 2.77 & 0.20 \\
% END VERIFIED R9.10 BY MINIMUM EPISODE LENGTH
\bottomrule
\end{tabular}\end{adjustbox}
\end{center}

\noindent The onset count, recall denominator and definition of a false alarm change across rows. The model, calibration, threshold, alarm stream and exposure do not.

\newpage
\subsection*{E. Threshold selection and denominator sensitivity}
\noindent The R9.9 threshold-selection role of development block 1 contains 36 runs (12 full simulated minutes). The code evaluated 19 cutoffs, $0.05$ through $0.95$ in steps of $0.05$, subject to at most two unmatched alerts per \emph{full} simulated minute. The feasible candidate with highest event $F_1$ was selected; ties were resolved by lower false-alert rate, higher recall and lower cutoff. Of 19 candidates, 17 satisfied the budget and 10 had a true alert. At cutoff 0.15, there were 22 matches, 24 false alerts and 159 observed events: $F_1=2(22)/(22+24+159)=0.2146$, false alarms/min $=24/12=2.00$. These statistics use the \emph{same selection role used to choose the cutoff}, and are not an independent validation of the selected operating point.

\medskip
\noindent For the five subsequent development evaluation blocks, the primary denominator is 200 observed frames per run at 0.1 seconds per frame. The secondary denominator is 0.1 seconds for each connected ($\beta_0=1$), decision-eligible scored frame at $t=4,\ldots,193$ in a run. Both rates use exactly the same unmatched alarm counts. Connected decision minutes are an exposure sensitivity; the numeric line of two/min was not independently specified for this second denominator. In Tables E--F, bold connected rates exceed that \emph{numerical reference}, not a separately established operating budget.

\begin{center}
\small
\setlength{\tabcolsep}{4pt}
\begin{adjustbox}{max width=\linewidth}\begin{tabular}{lrrrrrr}
\toprule
Mobility & False & Full min & False/full min & Connected frames & Connected min & False/connected min \\
\midrule
% BEGIN VERIFIED EXPOSURE SENSITIVITY POOLED
All runs & 88 & 60 & 1.47 & 28301 & 47.17 & 1.87 \\
Random waypoint (RWP) & 18 & 20 & 0.90 & 10130 & 16.88 & 1.07 \\
Gauss--Markov (GM) & 49 & 20 & 2.45 & 8472 & 14.12 & \textbf{3.47} \\
Mission-based & 21 & 20 & 1.05 & 9699 & 16.17 & 1.30 \\
% END VERIFIED EXPOSURE SENSITIVITY POOLED
\bottomrule
\end{tabular}\end{adjustbox}
\end{center}

\medskip
\noindent All runs contributed 34,200 scored decisions, of which 28,301 were connected. Full-exposure block-by-mobility budget exceedances: 6/15. Applying the same numerical two/min reference to connected decision exposure: 8/15. These are different denominator definitions applied to the \emph{same} 180 runs, not additional independent experiments.

\begin{center}
\small
\setlength{\tabcolsep}{3.8pt}
\begin{adjustbox}{max width=\linewidth}\begin{tabular}{llrrrrr}
\toprule
Block & Mobility & False & False/full min & Connected frames & Connected min & False/connected min \\
\midrule
% BEGIN VERIFIED EXPOSURE SENSITIVITY BLOCK MOBILITY
2 & Random waypoint (RWP) & 2 & 0.50 & 2139 & 3.56 & 0.56 \\
2 & Gauss--Markov (GM) & 10 & 2.50 & 1264 & 2.11 & \textbf{4.75} \\
2 & Mission-based & 7 & 1.75 & 1786 & 2.98 & \textbf{2.35} \\
3 & Random waypoint (RWP) & 0 & 0.00 & 2226 & 3.71 & 0.00 \\
3 & Gauss--Markov (GM) & 6 & 1.50 & 1678 & 2.80 & \textbf{2.15} \\
3 & Mission-based & 2 & 0.50 & 1765 & 2.94 & 0.68 \\
4 & Random waypoint (RWP) & 1 & 0.25 & 1784 & 2.97 & 0.34 \\
4 & Gauss--Markov (GM) & 9 & 2.25 & 1829 & 3.05 & \textbf{2.95} \\
4 & Mission-based & 0 & 0.00 & 2107 & 3.51 & 0.00 \\
5 & Random waypoint (RWP) & 9 & 2.25 & 2029 & 3.38 & \textbf{2.66} \\
5 & Gauss--Markov (GM) & 10 & 2.50 & 1786 & 2.98 & \textbf{3.36} \\
5 & Mission-based & 0 & 0.00 & 2075 & 3.46 & 0.00 \\
6 & Random waypoint (RWP) & 6 & 1.50 & 1952 & 3.25 & 1.84 \\
6 & Gauss--Markov (GM) & 14 & 3.50 & 1915 & 3.19 & \textbf{4.39} \\
6 & Mission-based & 12 & 3.00 & 1966 & 3.28 & \textbf{3.66} \\
% END VERIFIED EXPOSURE SENSITIVITY BLOCK MOBILITY
\bottomrule
\end{tabular}\end{adjustbox}
\end{center}

\noindent Two cells meet the full-exposure budget but exceed the two/min numerical reference when only connected decision time is counted: block 2 mission-based mobility and block 3 Gauss--Markov. Exposure arithmetic uses one eligible connected frame $=0.1$ seconds. This sensitivity does not reclassify the original prespecified full-exposure budget.


\newpage
\subsection*{F. Between-block variability and leave-one-block-out checks}
\noindent Each of the five evaluation blocks consists of nine distinct trajectory groups with four paired radio/graph runs per group (36 runs or 12 full simulated minutes). A trajectory-group seed is different across blocks, but that does not imply that shared simulation assumptions or the fitted model vary independently. A block is not an independent field trial. Counts and recall for all five blocks appear in Table 5 of the manuscript. Dropping one block at a time is an exploratory sensitivity check, not a held-out test and not an alternative inferential unit. The cutoff, calibration and alarm stream stay frozen.

\begin{center}
\small
\setlength{\tabcolsep}{5pt}
\begin{adjustbox}{max width=\linewidth}\begin{tabular}{rrrrr}
\toprule
Omitted block & Remaining groups & Remaining minutes & Overall false/min & GM false/min \\
\midrule
% BEGIN VERIFIED BLOCK LEAVE ONE OUT
2 & 36 & 48 & 1.44 & 2.44 \\
3 & 36 & 48 & 1.67 & 2.69 \\
4 & 36 & 48 & 1.62 & 2.50 \\
5 & 36 & 48 & 1.44 & 2.44 \\
6 & 36 & 48 & 1.17 & 2.19 \\
% END VERIFIED BLOCK LEAVE ONE OUT
\bottomrule
\end{tabular}\end{adjustbox}
\end{center}

\subsection*{G. Balanced trajectory-group resampling sensitivity}
\noindent Four radio/graph runs sharing a trajectory group were always sampled as one unit. We formed nine strata from three mobility types and three node counts, each with five groups (one per development block). For each of 12,000 deterministic, seeded replicates (seed 20261121), we resampled five groups with replacement within each stratum. Thus each replicate retains 45 trajectory groups, 180 paired simulation runs and 60 full simulated minutes, though some trajectory groups are repeated and others omitted. The table gives the empirical 2.5th and 97.5th percentiles. These are \emph{exploratory conditional resampling intervals}, not confirmatory population confidence intervals: the model and threshold remain fixed, all groups are development data, and possible shared block effects are not modeled. The resampling holds scenario proportions fixed; results do not justify comparison p-values, portability claims or a false-alarm guarantee.

\begin{center}
\small
\setlength{\tabcolsep}{4pt}
\begin{adjustbox}{max width=\linewidth}\begin{tabular}{lrrrrr}
\toprule
Scope & Groups & False/min & 2.5--97.5\% & Recall (\%) & 2.5--97.5\% recall \\
\midrule
% BEGIN VERIFIED CLUSTER BOOTSTRAP
All runs & 45 & 1.47 & [0.87, 2.10] & 16.4 & [9.6, 23.8] \\
Random waypoint & 15 & 0.90 & [0.15, 1.75] & 7.9 & [3.0, 14.6] \\
Gauss--Markov & 15 & 2.45 & [1.10, 3.85] & 25.2 & [12.9, 37.7] \\
Mission-based & 15 & 1.05 & [0.10, 2.15] & 12.2 & [0.8, 27.9] \\
% END VERIFIED CLUSTER BOOTSTRAP
\bottomrule
\end{tabular}\end{adjustbox}
\end{center}

\noindent A rate below two in one short block and above two in another establishes observed heterogeneity within this simulation design. It does not estimate the probability of future budget compliance for an independently deployed FANET.


\subsection*{H. Source-linked simulation and model details}

The frozen simulator configuration at \path{core/configs/r7_full_20seed_isolated.json} specifies a 2-D square of side 1,000~m; trajectory generation is implemented in \path{core/fanet/mobility.py} and \path{core/fanet/dataset.py}. Three mobility generators are replayed with 10, 20 and 30 nodes, nominal/degraded physical radio conditions and fixed/adaptive link caps. The configured physical-link model, not the alternative SINR model available elsewhere in the source tree, is used in these recorded runs.

The \path{dataset.adaptive_radius} function switches from 420~m to 567~m only if the simulated geometric density exceeds $2\times10^{-5}$. The \path{radio.received_power_dbm} and \path{radio.build_fixed_adaptive_adjacencies} functions implement the propagation abstraction and symmetric directed-link criterion. The model is the \path{models.FANETTopoGNN} class, with a two-layer GCN, geometric persistence image and learned gate; the \path{training.train_torch_model} branch for this model uses a single-snapshot dataset. The snapshot warm-up length of five restricts eligible decisions but is not an input-sequence length for the selected model. Frozen calibration and alarm selection appear in \path{frozen/r9_9_calibration_alert.py}.

\textbf{Evidence boundary:} these are source-verified simulation parameters and descriptive tests of a fixed decision stream, not a 3-D aerial mobility study, a packet-level simulator validation or a field radio study. The audit package provides the runnable code and frozen config; Section~H contains no new empirical results.

\newpage
\subsection*{I. Exploratory History-4 comparator: development-only comparison}
\noindent History-4 is a deterministic memory-of-fragmentation rule, $s(t)=\frac{1}{4}\sum_{k=1}^{4}\mathbf{1}[\beta_0(t-k)>1]$, using only the previous four observed network states. It is not a trained model, and $s(t)$ is not a calibrated probability. It uses the same connected-state alarm gate, 190 decision indices per run ($t=4,\ldots,193$), matching horizon, cooldown and full simulated-time denominator as the frozen TopoGNN evaluation. An independently regenerated threshold-selection block~1 (36 runs; 12 min) matched every $t+6$ target from the frozen R9.9 selection CSV. The 19-candidate threshold grid and event-$F_1$-subject-to-budget selection rule were the same for both systems, but their selected cutoffs were \emph{not} forced equal. The History-4 cutoff was $0.30$; its in-sample selection counts were 43 of 159 onsets matched, 21 false alarms, 1.75 false alarms/min and event $F_1=0.386$. The frozen TopoGNN selection cut-off remained $0.15$, with 22 of 159 onsets, 24 false alarms, 2.00/min and event $F_1=0.215$.

\begin{center}
\small
\setlength{\tabcolsep}{4pt}
\begin{adjustbox}{max width=\linewidth}\begin{tabular}{llrrrrr}
\toprule
Model & Stratum & Events & Matched & False & Event $F_1$ & False/min \\
\midrule
TopoGNN & All & 1114 & 183 & 88 & 0.264 & 1.47 \\
History-4 & All & 1114 & 381 & 135 & 0.467 & \textbf{2.25} \\
TopoGNN & RWP & 305 & 24 & 18 & 0.138 & 0.90 \\
History-4 & RWP & 305 & 102 & 49 & 0.447 & \textbf{2.45} \\
TopoGNN & GM & 465 & 117 & 49 & 0.371 & \textbf{2.45} \\
History-4 & GM & 465 & 158 & 43 & 0.474 & \textbf{2.15} \\
TopoGNN & Mission & 344 & 42 & 21 & 0.206 & 1.05 \\
History-4 & Mission & 344 & 121 & 43 & 0.476 & \textbf{2.15} \\
\bottomrule
\end{tabular}\end{adjustbox}
\end{center}

\noindent The same reference budget was used to \emph{select} operating points. The transferred History-4 operating point did not satisfy that budget on the 180 evaluation runs; a new cutoff was \emph{not} selected using evaluation labels. This post-hoc comparison was devised after inspection of the original R9.10 results, so these are NOT blind evaluation blocks for the new heuristic. It cannot support a comparative superiority claim. The 10,000-resample, mobility-by-node-count-stratified, trajectory-group paired descriptive percentile interval for TopoGNN minus History-4 event $F_1$ was $[-0.294,-0.118]$. That interval conditions on a retrospectively chosen heuristic and the frozen development sample; it is \emph{not} a confirmatory test or a basis for declaring model superiority. Auxiliary CSVs, the development selection sweep, source-code and audit receipt are available in \path{data_derived/wn_simple_baseline/}.

\newpage
\subsection*{J. Lead-time distribution and hypothetical action-latency sensitivity}
\noindent The event-level lead is the number of sampled frames from an emitted warning to its one-to-one matched fragmentation onset. The timestep is 0.1~s. The following analyses keep original alerts, match assignments, thresholds, event labels and false-alarm classifications \emph{fixed}. They do not run a controller or simulate prevention, packet routing, or actuation. Each assumed effective-action latency $L$ is the \emph{entire hypothetical interval between recorded alarm and achieved mitigation}; whether any actual system can meet that latency is unmeasured. A pair with lead $\geq L$ is only a \emph{potential temporal opportunity}, with boundary equality treated optimistically.

\begin{center}
\small
\setlength{\tabcolsep}{6pt}
\begin{adjustbox}{max width=\linewidth}\begin{tabular}{lrrrrrrr}
\toprule
Rule & 0.1 s & 0.2 s & 0.3 s & 0.4 s & 0.5 s & 0.6 s & Matched total \\
\midrule
% BEGIN VERIFIED LEAD HISTOGRAM
TopoGNN & 85 & 40 & 25 & 13 & 13 & 7 & 183 \\
History-4 & 170 & 79 & 40 & 40 & 24 & 28 & 381 \\
% END VERIFIED LEAD HISTOGRAM
\bottomrule
\end{tabular}\end{adjustbox}
\end{center}

\noindent The median among matched TopoGNN warnings was 0.2~s, but 85/183 (46.4\%) had only one tick of lead, and 125/183 (68.3\%) had at most two. History-4 is a retrospectively specified comparator with a different alert rate; these distributions cannot be interpreted as a fair controlled effectiveness comparison.

\begin{center}
\small
\setlength{\tabcolsep}{5pt}
\begin{adjustbox}{max width=\linewidth}\begin{tabular}{lrrrrrr}
\toprule
Rule & Assumed $L$ (s) & Matched & Potentially timely & \% matched & \% all onsets & False alerts \\
\midrule
% BEGIN VERIFIED ACTION LATENCY
TopoGNN & 0.1 & 183 & 183 & 100.0 & 16.4 & 88 \\
TopoGNN & 0.2 & 183 & 98 & 53.6 & 8.8 & 88 \\
TopoGNN & 0.3 & 183 & 58 & 31.7 & 5.2 & 88 \\
TopoGNN & 0.4 & 183 & 33 & 18.0 & 3.0 & 88 \\
TopoGNN & 0.5 & 183 & 20 & 10.9 & 1.8 & 88 \\
TopoGNN & 0.6 & 183 & 7 & 3.8 & 0.6 & 88 \\
History-4 & 0.1 & 381 & 381 & 100.0 & 34.2 & 135 \\
History-4 & 0.2 & 381 & 211 & 55.4 & 18.9 & 135 \\
History-4 & 0.3 & 381 & 132 & 34.6 & 11.8 & 135 \\
History-4 & 0.4 & 381 & 92 & 24.1 & 8.3 & 135 \\
History-4 & 0.5 & 381 & 52 & 13.6 & 4.7 & 135 \\
History-4 & 0.6 & 381 & 28 & 7.3 & 2.5 & 135 \\
% END VERIFIED ACTION LATENCY
\bottomrule
\end{tabular}\end{adjustbox}
\end{center}

\noindent Percentages in the penultimate column use all 1,114 observed onsets, \emph{not} only detected onsets. No 0.3-second action was measured or executed. History-4 generated 135 false alerts (2.25/full simulated minute), above the 2/min budget, so its higher count of potentially timely pairs does not imply that it is an acceptable operating point. The false-alert count is unchanged by the latency thought experiment, because late matched warnings are not reassigned as false alarms.

\begin{center}
\small
\setlength{\tabcolsep}{5pt}
\begin{adjustbox}{max width=\linewidth}\begin{tabular}{lrlrrr}
\toprule
Rule & $L$ (s) & Mobility & Total events & Matched & Potentially timely \\
\midrule
% BEGIN VERIFIED ACTION LATENCY BY MOBILITY
TopoGNN & 0.2 & RWP & 305 & 24 & 11 \\
TopoGNN & 0.2 & GM & 465 & 117 & 64 \\
TopoGNN & 0.2 & Mission & 344 & 42 & 23 \\
TopoGNN & 0.3 & RWP & 305 & 24 & 6 \\
TopoGNN & 0.3 & GM & 465 & 117 & 38 \\
TopoGNN & 0.3 & Mission & 344 & 42 & 14 \\
History-4 & 0.2 & RWP & 305 & 102 & 54 \\
History-4 & 0.2 & GM & 465 & 158 & 86 \\
History-4 & 0.2 & Mission & 344 & 121 & 71 \\
History-4 & 0.3 & RWP & 305 & 102 & 39 \\
History-4 & 0.3 & GM & 465 & 158 & 53 \\
History-4 & 0.3 & Mission & 344 & 121 & 40 \\
% END VERIFIED ACTION LATENCY BY MOBILITY
\bottomrule
\end{tabular}\end{adjustbox}
\end{center}

\noindent These subgroup rows reflect the same 45 trajectories and paired radio/graph variants as the main evaluation, with no independence or controller-action claims. Onset events may remain unpreventable even when there is enough lead time, and a late alert can still be operationally informative for a different task. The publication does not assess that question.


\newpage
\subsection*{K. Graph-connectivity model versus network-service evidence}
\noindent The following matrix distinguishes \emph{source code capability} from the \emph{mechanisms used in the R9.10 recorded evaluation}. Source paths are relative to the author-review package. Descriptive statements here do not add observations to the frozen 180-run evidence.

\begin{center}
\small
\setlength{\tabcolsep}{3.1pt}
\begin{adjustbox}{max width=\linewidth}\begin{tabular}{@{}p{0.21\linewidth}p{0.20\linewidth}p{0.49\linewidth}@{}}
\toprule
Layer or claim & R9.10 evidence & Scope and traceability \\
\midrule
Planar mobility & Included & 2-D positions/velocities; RWP, Gauss--Markov, mission. \path{core/fanet/mobility.py}; \path{core/fanet/dataset.py}. \\
Radio-proxy link graph & Included & \path{core/configs/r7_full_20seed_isolated.json}: \texttt{link\_model=physical}; received power with path loss, correlated shadowing, Rayleigh fading, receiver sensitivity and radius cap; both directed conditions required. \path{core/fanet/radio.py}. \\
SINR/interference link rule & Available in code, unused & \path{core/fanet/radio.py} defines an alternative \texttt{sinr} branch. The recorded R9.10 configuration selects \texttt{physical}, not \texttt{sinr}. \\
MAC, packet traffic and routing & Not measured & \path{core/fanet/packet_sim.py} contains a separate SimPy packet simulator, but R9.10 scores and ground truth were produced from topology snapshots, with no traffic or packet-log outputs. \\
Graph fragmentation onset & Direct outcome & \path{evaluation_full_truth.csv}: connected-component counts across frames. Matching evaluates thresholded-graph onset, not service outage. \\
PDR, throughput, end-to-end delay & Not measured & No packet arrivals, successful delivery counts, end-to-end timestamps or routing recovery measurements in the frozen R9.10 evaluation CSVs. \\
3-D flight, terrain and field test & Not performed & Neither altitude/terrain effects nor independent measurements calibrate the simulated link proxy to observed flight and network service. \\
\bottomrule
\end{tabular}\end{adjustbox}
\end{center}

\noindent The presence of a path-loss and fading model does \emph{not} imply a full PHY/MAC stack, and the presence of an optional packet simulator in a source archive does \emph{not} establish that it was executed in this experiment. Under this design a component split is an event in the thresholded link graph. Neither \emph{connected} nor \emph{fragmented} graph labels are empirically validated substitutes for a measured packet-delivery event. A future validation should (i) predefine packet loads and routing/MAC settings on matched trajectories, (ii) report graph-onset versus service-outage agreement, (iii) measure PDR, latency and routing recovery, and (iv) test warning-to-mitigation effects without tuning on evaluation traces. These steps are proposed follow-up work, not results of this paper.

\end{document}
